\subsection{交换因子}

在本段中，当我们不加说明地提及分次时，总是指类型为 \(\Delta\) 的分次，其中 \(\Delta\) 是一个以加法记写的交换群。在本段中， \(\Delta\) 上取值于 Z 的交换因子称为 \(\Delta\) 上的交换因子（ \(\S\) 4，第 7 号，定义 6）。因此， \(\Delta\) 上的一个交换因子等同于一个映射 \(\varepsilon\colon(\alpha,\beta)\mapsto\varepsilon_{\alpha\beta}=\varepsilon(\alpha,\beta)\) ，它把 \(\Delta\times\Delta\) 映到乘法群 \(\{-1,1\}\) 中，使得对 \(\alpha,\alpha^{\prime},\beta,\beta^{\prime}\) （均为 \(\Delta\) 中的元素），

\[
\left\{ \begin{array}{l} \varepsilon (\alpha + \alpha^ {\prime}, \beta) = \varepsilon (\alpha , \beta) \varepsilon (\alpha^ {\prime}, \beta) \\ \varepsilon (\alpha , \beta + \beta^ {\prime}) = \varepsilon (\alpha , \beta) \varepsilon (\alpha , \beta^ {\prime}) \\ \varepsilon (\beta , \alpha) = \varepsilon (\alpha , \beta). \end{array} \right.\tag{1}
\]

由此可得 \(\varepsilon(2\alpha, \beta) = \varepsilon(\alpha, 2\beta) = 1\) .

当 \(\Delta = \mathbf{Z}\) 时，每个交换因子 \(\varepsilon\) 由 \(\varepsilon(1, 1)\) 的值决定；因此这样的因子只有两个，第一个由

\[
\varepsilon (p, q) = 1 \quad \text { for } \quad p, q \text { in } \mathbf {Z}\tag{2}
\]

定义，第二个由

\[
\varepsilon (p, q) = (- 1) ^ {p q} \quad \text { for } \quad p, q \text { in } \mathbf {Z}.\tag{3}
\]

